\documentclass[10pt]{article}
\usepackage[utf8]{inputenc}
\usepackage[T1]{fontenc}
\usepackage[spanish,es-noquoting]{babel}
\usepackage[a4paper,landscape,includehead,top=0.8cm,bottom=0.8cm,left=1.2cm,right=1.2cm,headheight=14pt,headsep=5mm]{geometry}
\usepackage{longtable,array,xcolor,amssymb,amsmath,colortbl,tabularx,enumitem,fancyhdr,titlesec}
\usepackage{lastpage}

\definecolor{cab}{HTML}{1F3864}
\definecolor{gris}{HTML}{EFEFEF}
\definecolor{gris2}{HTML}{F8F8F8}
\definecolor{alerta}{HTML}{B00020}
\definecolor{clave}{HTML}{0B6E4F}
\definecolor{libre}{HTML}{DDDDDD}

\setlength{\parindent}{0pt}
\renewcommand{\arraystretch}{1.28}
\setlength{\tabcolsep}{3pt}

\pagestyle{fancy}\fancyhf{}
\lhead{\small\textbf{Cronograma --- Examen de admisión, Maestría en Física, Uniandes}}
\rhead{\small Examen: lunes 23 de noviembre de 2026 \quad|\quad Página \thepage\ de \pageref{LastPage}}
\renewcommand{\headrulewidth}{0.4pt}

\newcommand{\ch}{$\square$}
\newcommand{\seccion}[1]{\vspace{2mm}{\color{cab}\large\textbf{#1}}\vspace{1mm}\hrule\vspace{2mm}}

\titleformat{\section}{\color{cab}\Large\bfseries}{}{0pt}{}

\begin{document}


\begin{center}
{\color{cab}\Huge\textbf{Cronograma de estudio}}\\[2mm]
{\Large Examen de admisión --- Maestría en Ciencias, Física --- Universidad de los Andes}\\[3mm]
{\large\textbf{Del lunes 17 de agosto al domingo 22 de noviembre de 2026} \quad ---\quad 14 semanas, S0 a S13}\\[2mm]
{\large\color{alerta}\textbf{Examen: lunes 23 de noviembre de 2026}}
\end{center}

\vspace{1mm}
\small
\begin{minipage}[t]{0.48\textwidth}
\seccion{Cómo usar este documento}
Cada día tiene \textbf{dos versiones de la misma sesión}:
\begin{itemize}[leftmargin=5mm,itemsep=0.5mm]
\item \textbf{MÍNIMA (1 hora).} El piso. Si un día solo puedes esto, el plan sigue en pie.
\item \textbf{EXTENDIDA (2 horas).} La mínima \emph{más} lo que dice esa columna. No es una alternativa: es la mínima ampliada.
\end{itemize}
Marca la casilla \ch{} del día cuando termines. Si hiciste la extendida, marca también la de esa columna.

\seccion{Restricciones aplicadas}
\begin{itemize}[leftmargin=5mm,itemsep=0.5mm]
\item Lunes, martes y miércoles: desde las \textbf{16:00}.
\item Jueves y viernes: desde las \textbf{19:00}.
\item \textbf{El REA (Molitoris) acompaña cada semana.} Su repaso son 74 páginas para las siete áreas, escritas para este formato de opción múltiple. Se lee antes del texto grande, no en vez de él. Sus cuatro exámenes con explicaciones detalladas son la munición de las semanas 12 y 13.
\item \textbf{Sábados: comodín.} Sin asignación fija; se usan solo si quedó algo pendiente entre semana.
\item \textbf{Domingos: comodín.} No traen material nuevo. Son para recuperar, cerrar el entregable de la semana o descansar. Puedes dejarlos libres sin romper nada.
\item \textbf{Los días 12 de cada mes están libres} (12 de septiembre, 12 de octubre y 12 de noviembre). Si caen entre semana, la tarea de ese día se corre al domingo comodín.
\end{itemize}

\seccion{Presupuesto de horas --- léelo}
Cinco sesiones semanales de 1 a 2 horas dan entre \textbf{5 y 10 horas por semana}, es decir entre 70 y 140 horas en total. Es un presupuesto ajustado para este examen.\\[1mm]
Consecuencia práctica: \textbf{el plan está podado al hueso}. No hay lectura de relleno ni capítulos «por cultura». Si un día tienes que elegir, haz siempre la MÍNIMA completa antes que media EXTENDIDA.\\[1mm]
Si puedes sostener la EXTENDIDA tres o más días por semana, llegas holgado. Si solo alcanzas las mínimas, llegas justo pero llegas.
\end{minipage}
\hfill
\begin{minipage}[t]{0.48\textwidth}
\seccion{Dónde está cada cosa}
Todo cuelga de \texttt{\textasciitilde/Documents/ANDES/}:
\begin{itemize}[leftmargin=5mm,itemsep=0.3mm]
\item \texttt{02\_BIBLIOTECA/} --- los libros, ordenados por área.
\item \texttt{01\_EXAMENES/uniandes\_admision/} --- el simulacro oficial.
\item \texttt{01\_EXAMENES/gre\_physics/} --- 4 exámenes GRE con soluciones, ETS y el REA.
\item \texttt{01\_EXAMENES/banco\_hrw/} --- 2.200 preguntas con respuesta.
\item \texttt{01\_EXAMENES/otros\_bancos/mit\_ocw/} --- problem sets y exámenes del MIT.
\item \texttt{03\_TEMARIO/} --- el desglose detallado de cada semana.
\item \texttt{04\_ESTUDIO/formularios/} --- aquí escribes tus formularios.
\item \texttt{06\_SEGUIMIENTO/} --- el medidor. Lee su \texttt{INSTRUCTIVO.md}.
\end{itemize}

\seccion{Las tres reglas}
\begin{enumerate}[leftmargin=5mm,itemsep=0.8mm]
\item \textbf{El viernes es sagrado.} Siempre opción múltiple cronometrada. Es el único entrenamiento del formato real. Sin esto, sabrás física pero no aprobarás.
\item \textbf{Todo fallo se anota.} Con \texttt{medidor.py error}, o al registrar el simulacro, indicando la causa: concepto, álgebra, lectura, tiempo o descuido. Se rehace a los 3 y a los 14 días.
\item \textbf{Los 10 patrones mandan.} Están marcados en negrita a lo largo del cronograma. Son el 40\,\% del examen y lo que separa a quien entra de quien no.
\end{enumerate}

\seccion{Ritmo objetivo por problema}
\begin{tabular}{@{}ll@{}}
Semanas 1--3 & Resolver bien, sin reloj \\
Semanas 4--9 & Con reloj: 10 min por problema \\
Semanas 10--12 & 7 min por problema \\
Semanas 13--14 & 6 min, en simulacro completo \\
\end{tabular}\\[2mm]
En el examen real tienes \textbf{7,2 minutos por pregunta}: 25 preguntas en 3 horas.
\end{minipage}

\newpage

\seccion{Mapa de habilidades}
Una fila por habilidad evaluada en el examen. Cada semana escribe tu nivel del 0 al 4. Misma escala que
\texttt{medidor.py skills}, para contrastar lo que sientes con lo que miden los simulacros.\\[0.5mm]
\textbf{0} no lo he visto \quad \textbf{1} lo vi, no me sale \quad \textbf{2} lo entiendo si lo miro \quad
\textbf{3} lo resuelvo solo \quad \textbf{4} lo resuelvo rápido\\[1mm]
Las filas en negrita son los \textbf{10 patrones}: son el 40\,\% del examen. Ninguna debería quedar bajo 3 en noviembre.
\par\vspace{2.5mm}
{\renewcommand{\arraystretch}{1.34}\setlength{\tabcolsep}{5.2pt}
\noindent\begin{tabular}{@{}|p{10.4cm}|c|c|c|c|c|c|c|c|c|c|c|c|c|c|}
\hline
\rowcolor{cab}\color{white}\textbf{\ Habilidad} & \color{white}\scriptsize\textbf{S1} & \color{white}\scriptsize\textbf{S2} & \color{white}\scriptsize\textbf{S3} & \color{white}\scriptsize\textbf{S4} & \color{white}\scriptsize\textbf{S5} & \color{white}\scriptsize\textbf{S6} & \color{white}\scriptsize\textbf{S7} & \color{white}\scriptsize\textbf{S8} & \color{white}\scriptsize\textbf{S9} & \color{white}\scriptsize\textbf{S10} & \color{white}\scriptsize\textbf{S11} & \color{white}\scriptsize\textbf{S12} & \color{white}\scriptsize\textbf{S13} & \color{white}\scriptsize\textbf{S14}\\ \hline
\rowcolor{gris}\multicolumn{15}{|l|}{\textbf{\small Mecánica}}\\ \hline
\scriptsize Conservación de energía + condición de contacto en rizo vertical &   &   &   &   &   &   &   &   &   &   &   &   &   &  \\ \hline
\scriptsize Newton con ligaduras y fricción estática (Sears 5.88) &   &   &   &   &   &   &   &   &   &   &   &   &   &  \\ \hline
\scriptsize Cinemática rotacional con aceleración angular constante (Sears 9.15) &   &   &   &   &   &   &   &   &   &   &   &   &   &  \\ \hline
\scriptsize Energía potencial elástica &   &   &   &   &   &   &   &   &   &   &   &   &   &  \\ \hline
\scriptsize Oscilador forzado amortiguado: amplitud en estado estacionario &   &   &   &   &   &   &   &   &   &   &   &   &   &  \\ \hline
\scriptsize \textbf{P1 — Área en espacio de fase y período T(E)=dA/dE} &   &   &   &   &   &   &   &   &   &   &   &   &   &  \\ \hline
\scriptsize \textbf{P2 — Corchetes de Poisson} &   &   &   &   &   &   &   &   &   &   &   &   &   &  \\ \hline
\scriptsize \textbf{P3 — Scattering con esfera dura: parámetro de impacto b(theta)} &   &   &   &   &   &   &   &   &   &   &   &   &   &  \\ \hline
\rowcolor{gris}\multicolumn{15}{|l|}{\textbf{\small Electromagnetismo}}\\ \hline
\scriptsize Ley de Ampère en conductor cilíndrico hueco &   &   &   &   &   &   &   &   &   &   &   &   &   &  \\ \hline
\scriptsize \textbf{P4 — Condición de gauge en potenciales EM} &   &   &   &   &   &   &   &   &   &   &   &   &   &  \\ \hline
\scriptsize \textbf{P5 — Esfera conductora a tierra en campo uniforme (armónicos esféricos)} &   &   &   &   &   &   &   &   &   &   &   &   &   &  \\ \hline
\scriptsize Campo B en el centro de una espira circular &   &   &   &   &   &   &   &   &   &   &   &   &   &  \\ \hline
\scriptsize Superposición vectorial de campos de dos anillos cargados &   &   &   &   &   &   &   &   &   &   &   &   &   &  \\ \hline
\scriptsize Diferencia de potencial entre esfera sólida y cascarón &   &   &   &   &   &   &   &   &   &   &   &   &   &  \\ \hline
\rowcolor{gris}\multicolumn{15}{|l|}{\textbf{\small Relatividad}}\\ \hline
\scriptsize Dilatación temporal con expansión para beta << 1 &   &   &   &   &   &   &   &   &   &   &   &   &   &  \\ \hline
\scriptsize Doppler relativista y condición de ionización &   &   &   &   &   &   &   &   &   &   &   &   &   &  \\ \hline
\rowcolor{gris}\multicolumn{15}{|l|}{\textbf{\small Termo y estadística}}\\ \hline
\scriptsize Coeficiente de expansión lineal a partir de L(T) &   &   &   &   &   &   &   &   &   &   &   &   &   &  \\ \hline
\scriptsize Calorimetría con dos cambios de fase &   &   &   &   &   &   &   &   &   &   &   &   &   &  \\ \hline
\scriptsize Primera ley con trabajo sobre trayectoria p(V) no trivial &   &   &   &   &   &   &   &   &   &   &   &   &   &  \\ \hline
\scriptsize \textbf{P6 — Entropía combinatoria S = k ln Omega} &   &   &   &   &   &   &   &   &   &   &   &   &   &  \\ \hline
\scriptsize \textbf{P7 — Dos sistemas en contacto térmico: energía más probable} &   &   &   &   &   &   &   &   &   &   &   &   &   &  \\ \hline
\rowcolor{gris}\multicolumn{15}{|l|}{\textbf{\small Moderna y cuántica}}\\ \hline
\scriptsize \textbf{P8 — Dos bosones indistinguibles en caja rectangular 2D} &   &   &   &   &   &   &   &   &   &   &   &   &   &  \\ \hline
\scriptsize \textbf{P9 — Escalón de potencial: coeficiente de reflexión R} &   &   &   &   &   &   &   &   &   &   &   &   &   &  \\ \hline
\scriptsize \textbf{P10 — Oscilador armónico cuántico: varianza como función del tiempo} &   &   &   &   &   &   &   &   &   &   &   &   &   &  \\ \hline
\end{tabular}}
\newpage
\seccion{S0 \quad\textnormal{\normalsize 17 de agosto -- 23 de agosto}}
\textbf{Foco:} Diagnóstico y los 10 patrones\\[0.5mm]
\textbf{Por qué:} El diagnóstico ya dio su número: 2 de 12, con la base en 0/6 y los patrones en 2/6. Esta semana se cierra la otra mitad y se convierte el resultado en los 10 patrones escritos a mano.\\[0.5mm]
\textbf{Entregable de la semana:} \texttt{04\_ESTUDIO/formularios/patrones.md} con los 10 derivados, no copiados.
\par\vspace{2.5mm}
\noindent\begin{tabular}{@{}p{0.45cm}p{2.5cm}p{1.05cm}|p{0.45cm}p{9.75cm}|p{0.45cm}p{9.75cm}@{}}
\rowcolor{cab}
\multicolumn{3}{@{}l}{\color{white}\textbf{\ Día}} & \multicolumn{2}{l}{\color{white}\textbf{MÍNIMA --- 1 hora}} & \multicolumn{2}{l}{\color{white}\textbf{EXTENDIDA --- 2 horas (la mínima + esto)}}\\
\ch & \textbf{Lunes 17} & 16:00 & \ch & Los 10 patrones de 03\_TEMARIO/00\_mapa\_simulacro.md. Empieza por los cuatro que fallaste: P2 Poisson, P7 contacto térmico, P8 bosones en caja, P9 escalón. & \ch & Además: P3, P4, P6 y P10, que el diagnóstico no llegó a evaluar.\\
\rowcolor{gris2}
\ch & \textbf{Martes 18} & 16:00 & \ch & Cierra el diagnóstico: preguntas 13--24 del simulacro, 60 min cronometrados, sin apuntes. & \ch & Además: regístralo con \texttt{python3 06\_SEGUIMIENTO/medidor.py simulacro}.\\
\ch & \textbf{Miércoles 19} & 16:00 & \ch & Corre \texttt{medidor.py reporte} y abre \texttt{tablero.html}. Identifica tus 2 áreas más débiles. & \ch & Además: al registrar el simulacro, clasifica cada fallo por causa. Es el paso que más informa.\\
\rowcolor{gris2}
\ch & \textbf{Jueves 20} & 19:00 & \ch & Sears Vol.\,1, cinemática 1D y 2D: lee el resumen del capítulo y resuelve 8 problemas de fin de capítulo. & \ch & Además: 8 problemas más + Irodov §1.1 (p.\,11), 5 problemas.\\
\ch & \textbf{Viernes 21} & 19:00 & \ch & Banco HRW caps.\,2--4: 30 preguntas cronometradas a 2 min. Corrige con el «Ans:» al pie de cada una. & \ch & 60 preguntas en lugar de 30.\\
\rowcolor{libre}
 & \textbf{Sábado 22} & --- & \multicolumn{4}{l@{}}{\textit{Comodín. Sin asignación fija.}}\\
\rowcolor{gris}
\ch & \textbf{Domingo 23} & libre & \multicolumn{4}{l@{}}{\textit{Comodín. Recupera lo que no alcanzaste. Si vas al día, rehaz los 5 fallos más graves del diagnóstico.}}\\
\end{tabular}
\vspace{2mm}

\noindent\begin{tabular}{@{}p{2.0cm}p{22.8cm}@{}}
\textbf{\small Notas} & \rule{\linewidth}{0.4pt}\\[5mm]
 & \rule{\linewidth}{0.4pt}\\[5mm]
 & \rule{\linewidth}{0.4pt}\\[5mm]
 & \rule{\linewidth}{0.4pt}\\[5mm]
 & \rule{\linewidth}{0.4pt}\\
\end{tabular}
\newpage
\seccion{S1 \quad\textnormal{\normalsize 24 de agosto -- 30 de agosto}}
\textbf{Foco:} Newton, energía y momento\\[0.5mm]
\textbf{Por qué:} Consolidar el núcleo de Física 1. Aquí viven 4 de las 24 preguntas del simulacro.\\[0.5mm]
\textbf{Entregable de la semana:} 04\_ESTUDIO/formularios/mecanica\_I.md
\par\vspace{2.5mm}
\noindent\begin{tabular}{@{}p{0.45cm}p{2.5cm}p{1.05cm}|p{0.45cm}p{9.75cm}|p{0.45cm}p{9.75cm}@{}}
\rowcolor{cab}
\multicolumn{3}{@{}l}{\color{white}\textbf{\ Día}} & \multicolumn{2}{l}{\color{white}\textbf{MÍNIMA --- 1 hora}} & \multicolumn{2}{l}{\color{white}\textbf{EXTENDIDA --- 2 horas (la mínima + esto)}}\\
\ch & \textbf{Lunes 24} & 16:00 & \ch & Sears Vol.\,1, leyes de Newton y fricción: resumen + 8 problemas. Incluye el \textbf{problema 5.88} (es la pregunta 2 del examen real). & \ch & Además: Morin cap.\,3 (p.\,51), lee 3 ejemplos resueltos y haz 5 problemas.\\
\rowcolor{gris2}
\ch & \textbf{Martes 25} & 16:00 & \ch & Sears, trabajo y energía: resumen + 8 problemas. & \ch & Además: Morin cap.\,5 (p.\,138), 6 problemas con solución completa.\\
\ch & \textbf{Miércoles 26} & 16:00 & \ch & Sears, momento lineal, impulso y colisiones: 8 problemas. & \ch & Además: Irodov §1.3 (p.\,30), 6 problemas.\\
\rowcolor{gris2}
\ch & \textbf{Jueves 27} & 19:00 & \ch & Repaso activo: rehaz \emph{sin mirar} 5 problemas que fallaste esta semana. & \ch & Además: Morin cap.\,2 estática (p.\,22), 5 problemas.\\
\ch & \textbf{Viernes 28} & 19:00 & \ch & Banco HRW caps.\,5--9: 30 preguntas a 2 min. & \ch & 60 preguntas.\\
\rowcolor{libre}
 & \textbf{Sábado 29} & --- & \multicolumn{4}{l@{}}{\textit{Comodín. Sin asignación fija.}}\\
\rowcolor{gris}
\ch & \textbf{Domingo 30} & libre & \multicolumn{4}{l@{}}{\textit{Comodín. Cierra el formulario de mecánica I si quedó pendiente.}}\\
\end{tabular}
\vspace{2mm}

\noindent\begin{tabular}{@{}p{2.0cm}p{22.8cm}@{}}
\textbf{\small Notas} & \rule{\linewidth}{0.4pt}\\[5mm]
 & \rule{\linewidth}{0.4pt}\\[5mm]
 & \rule{\linewidth}{0.4pt}\\[5mm]
 & \rule{\linewidth}{0.4pt}\\[5mm]
 & \rule{\linewidth}{0.4pt}\\
\end{tabular}
\newpage
\seccion{S2 \quad\textnormal{\normalsize 31 de agosto -- 6 de septiembre}}
\textbf{Foco:} Rotación, gravitación, oscilaciones y fluidos\\[0.5mm]
\textbf{Por qué:} Terminar Física 1. La fórmula del oscilador forzado es pregunta directa del examen.\\[0.5mm]
\textbf{Entregable de la semana:} 04\_ESTUDIO/formularios/mecanica\_II.md + 1 tarjeta Anki (amplitud del oscilador forzado).
\par\vspace{2.5mm}
\noindent\begin{tabular}{@{}p{0.45cm}p{2.5cm}p{1.05cm}|p{0.45cm}p{9.75cm}|p{0.45cm}p{9.75cm}@{}}
\rowcolor{cab}
\multicolumn{3}{@{}l}{\color{white}\textbf{\ Día}} & \multicolumn{2}{l}{\color{white}\textbf{MÍNIMA --- 1 hora}} & \multicolumn{2}{l}{\color{white}\textbf{EXTENDIDA --- 2 horas (la mínima + esto)}}\\
\ch & \textbf{Lunes 31} & 16:00 & \ch & Sears, cinemática y dinámica rotacional, momento de inercia: 8 problemas. Incluye el \textbf{problema 9.15} (es la pregunta 3 del examen real). & \ch & Además: Morin cap.\,8 (p.\,309), 5 problemas.\\
\rowcolor{gris2}
\ch & \textbf{Martes 1} & 16:00 & \ch & Sears, momento angular, rodadura y torque: 8 problemas. & \ch & Además: Irodov §1.5 (p.\,47), 6 problemas.\\
\ch & \textbf{Miércoles 2} & 16:00 & \ch & Sears, gravitación y leyes de Kepler: 6 problemas. Luego fluidos (Arquímedes, Bernoulli): 4 problemas. & \ch & Además: Irodov §1.4 (p.\,43) y §1.7 (p.\,62), 8 problemas.\\
\rowcolor{gris2}
\ch & \textbf{Jueves 3} & 19:00 & \ch & Oscilaciones. Morin cap.\,4 (p.\,101), 6 problemas. \textbf{Clave:} oscilador forzado amortiguado --- memoriza la amplitud en estado estacionario. & \ch & Además: MIT 8.03 ProblemSet1 a ProblemSet3.\\
\ch & \textbf{Viernes 4} & 19:00 & \ch & Banco HRW caps.\,10--15: 30 preguntas a 2 min. & \ch & 60 preguntas.\\
\rowcolor{libre}
 & \textbf{Sábado 5} & --- & \multicolumn{4}{l@{}}{\textit{Comodín. Sin asignación fija.}}\\
\rowcolor{gris}
\ch & \textbf{Domingo 6} & libre & \multicolumn{4}{l@{}}{\textit{Comodín. Escribe el formulario de mecánica II y la tarjeta del oscilador forzado.}}\\
\end{tabular}
\vspace{2mm}

\noindent\begin{tabular}{@{}p{2.0cm}p{22.8cm}@{}}
\textbf{\small Notas} & \rule{\linewidth}{0.4pt}\\[5mm]
 & \rule{\linewidth}{0.4pt}\\[5mm]
 & \rule{\linewidth}{0.4pt}\\[5mm]
 & \rule{\linewidth}{0.4pt}\\[5mm]
 & \rule{\linewidth}{0.4pt}\\
\end{tabular}
\newpage
\seccion{S3 \quad\textnormal{\normalsize 7 de septiembre -- 13 de septiembre}}
\textbf{Foco:} Mecánica analítica --- LA SEMANA MÁS RENTABLE\\[0.5mm]
\textbf{Por qué:} Tres preguntas del examen (6, 7 y 8) salen de aquí y son inaccesibles sin haberlas visto. 12\,\% del examen en 5 sesiones.\\[0.5mm]
\textbf{Entregable de la semana:} 04\_ESTUDIO/formularios/mecanica\_analitica.md con los 3 patrones + 3 tarjetas Anki. \textbf{No pases a la semana 5 sin esto.}
\par\vspace{2.5mm}
\noindent\begin{tabular}{@{}p{0.45cm}p{2.5cm}p{1.05cm}|p{0.45cm}p{9.75cm}|p{0.45cm}p{9.75cm}@{}}
\rowcolor{cab}
\multicolumn{3}{@{}l}{\color{white}\textbf{\ Día}} & \multicolumn{2}{l}{\color{white}\textbf{MÍNIMA --- 1 hora}} & \multicolumn{2}{l}{\color{white}\textbf{EXTENDIDA --- 2 horas (la mínima + esto)}}\\
\ch & \textbf{Lunes 7} & 16:00 & \ch & Tong \emph{Classical Dynamics}: coordenadas generalizadas y ecuaciones de Lagrange. Lee y \textbf{rededuce a mano}. + Morin cap.\,6 (p.\,218), 3 problemas. & \ch & Además: 5 problemas más de Morin cap.\,6.\\
\rowcolor{gris2}
\ch & \textbf{Martes 8} & 16:00 & \ch & Tong: hamiltoniano, transformación de Legendre, ecuaciones canónicas. Rededuce. + Lim \emph{Mechanics} 2068--2084, 3 problemas. & \ch & Además: 4 problemas más de ese rango.\\
\ch & \textbf{Miércoles 9} & 16:00 & \ch & \textbf{PATRÓN 1 --- Espacio de fase:} $T(E)=dA/dE$. Entiende la derivación completa y escríbela de memoria. Tong, variables acción-ángulo. & \ch & Además: Lim \emph{Mechanics} 2028--2067 (pequeñas oscilaciones), 4 problemas.\\
\rowcolor{gris2}
\ch & \textbf{Jueves 10} & 19:00 & \ch & \textbf{PATRÓN 2 --- Corchetes de Poisson:} $\{q_i,p_j\}=\delta_{ij}$, $\{L_i,p_j\}=\varepsilon_{ijk}p_k$, $\{L_i,L_j\}=\varepsilon_{ijk}L_k$. Derívalos tú, no los copies. & \ch & Además: Goldstein, sección de corchetes de Poisson: 4 ejercicios.\\
\ch & \textbf{Viernes 11} & 19:00 & \ch & \textbf{PATRÓN 3 --- Scattering esfera dura:} $b=R\cos(\theta/2)$, $d\sigma/d\Omega=R^2/4$, $\sigma=\pi R^2$. Deriva y memoriza. Luego: preguntas de mecánica analítica de GR8677 y GR9277. & \ch & Además: Lim \emph{Mechanics} 2001--2027, 5 problemas.\\
\rowcolor{libre}
 & \textbf{Sábado 12} & --- & \multicolumn{4}{l@{}}{\textit{Comodín. Sin asignación fija.}}\\
\rowcolor{gris}
\ch & \textbf{Domingo 13} & libre & \multicolumn{4}{l@{}}{\textit{Comodín. Si algún patrón no quedó, hoy es el día. Es la semana que no se puede dejar a medias.}}\\
\end{tabular}
\vspace{2mm}

\noindent\begin{tabular}{@{}p{2.0cm}p{22.8cm}@{}}
\textbf{\small Notas} & \rule{\linewidth}{0.4pt}\\[5mm]
 & \rule{\linewidth}{0.4pt}\\[5mm]
 & \rule{\linewidth}{0.4pt}\\[5mm]
 & \rule{\linewidth}{0.4pt}\\[5mm]
 & \rule{\linewidth}{0.4pt}\\
\end{tabular}
\newpage
\seccion{S4 \quad\textnormal{\normalsize 14 de septiembre -- 20 de septiembre}}
\textbf{Foco:} Electrostática\\[0.5mm]
\textbf{Por qué:} Abrir el bloque de E\&M, que pesa 28\,\%. Griffiths es el texto de las próximas tres semanas.\\[0.5mm]
\textbf{Entregable de la semana:} 04\_ESTUDIO/formularios/em\_I.md
\par\vspace{2.5mm}
\noindent\begin{tabular}{@{}p{0.45cm}p{2.5cm}p{1.05cm}|p{0.45cm}p{9.75cm}|p{0.45cm}p{9.75cm}@{}}
\rowcolor{cab}
\multicolumn{3}{@{}l}{\color{white}\textbf{\ Día}} & \multicolumn{2}{l}{\color{white}\textbf{MÍNIMA --- 1 hora}} & \multicolumn{2}{l}{\color{white}\textbf{EXTENDIDA --- 2 horas (la mínima + esto)}}\\
\ch & \textbf{Lunes 14} & 16:00 & \ch & Griffiths cap.\,2: Coulomb, campo eléctrico, ley de Gauss (§2.1--2.2). 4 problemas. & \ch & Además: 4 problemas más + Irodov §3.1 (p.\,105), 5 problemas.\\
\rowcolor{gris2}
\ch & \textbf{Martes 15} & 16:00 & \ch & Griffiths §2.3, potencial eléctrico --- de ahí sale la \textbf{pregunta 15}. 5 problemas. & \ch & Además: Sears Vol.\,2, potencial: 8 problemas.\\
\ch & \textbf{Miércoles 16} & 16:00 & \ch & Conductores y capacitancia: Griffiths §2.5 + Sears Vol.\,2. 8 problemas. & \ch & Además: Irodov §3.3 (p.\,118), 6 problemas.\\
\rowcolor{gris2}
\ch & \textbf{Jueves 17} & 19:00 & \ch & Dieléctricos: Griffiths cap.\,4 (p.\,160), §4.1--4.3. 4 problemas. & \ch & Además: Lim \emph{Electromagnetism} 1043--1061, 4 problemas.\\
\ch & \textbf{Viernes 18} & 19:00 & \ch & Banco HRW caps.\,21--25: 30 preguntas a 2 min. & \ch & 60 preguntas.\\
\rowcolor{libre}
 & \textbf{Sábado 19} & --- & \multicolumn{4}{l@{}}{\textit{Comodín. Sin asignación fija.}}\\
\rowcolor{gris}
\ch & \textbf{Domingo 20} & libre & \multicolumn{4}{l@{}}{\textit{Comodín. Formulario de electrostática.}}\\
\end{tabular}
\vspace{2mm}

\noindent\begin{tabular}{@{}p{2.0cm}p{22.8cm}@{}}
\textbf{\small Notas} & \rule{\linewidth}{0.4pt}\\[5mm]
 & \rule{\linewidth}{0.4pt}\\[5mm]
 & \rule{\linewidth}{0.4pt}\\[5mm]
 & \rule{\linewidth}{0.4pt}\\[5mm]
 & \rule{\linewidth}{0.4pt}\\
\end{tabular}
\newpage
\seccion{S5 \quad\textnormal{\normalsize 21 de septiembre -- 27 de septiembre}}
\textbf{Foco:} Magnetostática, circuitos e inducción\\[0.5mm]
\textbf{Por qué:} La mitad mecánica del E\&M. Los circuitos son preguntas rápidas: aprovéchalas.\\[0.5mm]
\textbf{Entregable de la semana:} 04\_ESTUDIO/formularios/em\_II.md
\par\vspace{2.5mm}
\noindent\begin{tabular}{@{}p{0.45cm}p{2.5cm}p{1.05cm}|p{0.45cm}p{9.75cm}|p{0.45cm}p{9.75cm}@{}}
\rowcolor{cab}
\multicolumn{3}{@{}l}{\color{white}\textbf{\ Día}} & \multicolumn{2}{l}{\color{white}\textbf{MÍNIMA --- 1 hora}} & \multicolumn{2}{l}{\color{white}\textbf{EXTENDIDA --- 2 horas (la mínima + esto)}}\\
\ch & \textbf{Lunes 21} & 16:00 & \ch & Griffiths cap.\,5: fuerza de Lorentz y Biot--Savart (§5.2, p.\,215). Memoriza $B=\mu_0 I/2r$ en el centro de una espira (\textbf{pregunta 13}). 5 problemas. & \ch & Además: Lim \emph{EM} 2001--2038, 5 problemas.\\
\rowcolor{gris2}
\ch & \textbf{Martes 22} & 16:00 & \ch & Griffiths §5.3.3, aplicaciones de Ampère (p.\,225) --- de ahí sale la \textbf{pregunta 9}. Resuelve el conductor cilíndrico hueco en las 3 regiones. & \ch & Además: Irodov §3.5 (p.\,136), 6 problemas.\\
\ch & \textbf{Miércoles 23} & 16:00 & \ch & Corriente, resistencia y circuitos: Sears Vol.\,2 + Lim \emph{EM} 3001--3026. 10 problemas. & \ch & Además: Irodov §3.4 (p.\,125), 8 problemas.\\
\rowcolor{gris2}
\ch & \textbf{Jueves 24} & 19:00 & \ch & Faraday e inductancia: Griffiths §7.1--7.2 (p.\,285--320). 5 problemas. & \ch & Además: Lim \emph{EM} 2039--2063, 5 problemas.\\
\ch & \textbf{Viernes 25} & 19:00 & \ch & Banco HRW caps.\,26--31: 30 preguntas a 2 min. & \ch & 60 preguntas.\\
\rowcolor{libre}
 & \textbf{Sábado 26} & --- & \multicolumn{4}{l@{}}{\textit{Comodín. Sin asignación fija.}}\\
\rowcolor{gris}
\ch & \textbf{Domingo 27} & libre & \multicolumn{4}{l@{}}{\textit{Comodín. Formulario de magnetostática e inducción.}}\\
\end{tabular}
\vspace{2mm}

\noindent\begin{tabular}{@{}p{2.0cm}p{22.8cm}@{}}
\textbf{\small Notas} & \rule{\linewidth}{0.4pt}\\[5mm]
 & \rule{\linewidth}{0.4pt}\\[5mm]
 & \rule{\linewidth}{0.4pt}\\[5mm]
 & \rule{\linewidth}{0.4pt}\\[5mm]
 & \rule{\linewidth}{0.4pt}\\
\end{tabular}
\newpage
\seccion{S6 \quad\textnormal{\normalsize 28 de septiembre -- 4 de octubre}}
\textbf{Foco:} Maxwell, gauge y ondas EM --- SEGUNDA SEMANA MÁS RENTABLE\\[0.5mm]
\textbf{Por qué:} Aquí están las preguntas 10 y 12, los dos patrones caros de E\&M.\\[0.5mm]
\textbf{Entregable de la semana:} 04\_ESTUDIO/formularios/em\_III.md + 2 tarjetas Anki (gauge, armónicos esféricos).
\par\vspace{2.5mm}
\noindent\begin{tabular}{@{}p{0.45cm}p{2.5cm}p{1.05cm}|p{0.45cm}p{9.75cm}|p{0.45cm}p{9.75cm}@{}}
\rowcolor{cab}
\multicolumn{3}{@{}l}{\color{white}\textbf{\ Día}} & \multicolumn{2}{l}{\color{white}\textbf{MÍNIMA --- 1 hora}} & \multicolumn{2}{l}{\color{white}\textbf{EXTENDIDA --- 2 horas (la mínima + esto)}}\\
\ch & \textbf{Lunes 28} & 16:00 & \ch & Griffiths §7.3 (p.\,321): ecuaciones de Maxwell completas y corriente de desplazamiento. Escríbelas de memoria. & \ch & Además: cap.\,8 (p.\,345), vector de Poynting.\\
\rowcolor{gris2}
\ch & \textbf{Martes 29} & 16:00 & \ch & \textbf{PATRÓN 4 --- Gauge.} Griffiths §10.1.2 (p.\,419) y §10.1.3 (p.\,421). De ahí sale la \textbf{pregunta 10}. Deriva la condición de Lorenz y la de Coulomb. & \ch & Además: Tong \emph{Electromagnetism}, sección de potenciales.\\
\ch & \textbf{Miércoles 30} & 16:00 & \ch & \textbf{PATRÓN 5 --- Armónicos esféricos.} Griffiths §3.3.2 (p.\,137): esfera conductora a tierra en campo uniforme. Es la \textbf{pregunta 12}. Resuélvela entera. & \ch & Además: Griffiths §3.2, método de imágenes (p.\,121): 4 problemas.\\
\rowcolor{gris2}
\ch & \textbf{Jueves 1} & 19:00 & \ch & Ondas EM: Griffiths §9.1--9.2 (p.\,364--381). Polarización, reflexión y transmisión. 4 problemas. & \ch & Además: Lim \emph{EM} 1062--1095, 5 problemas.\\
\ch & \textbf{Viernes 2} & 19:00 & \ch & Preguntas de ondas EM y potenciales de GR9677 y GR0177, cronometradas. & \ch & Además: Banco HRW caps.\,32--33, 30 preguntas.\\
\rowcolor{libre}
 & \textbf{Sábado 3} & --- & \multicolumn{4}{l@{}}{\textit{Comodín. Sin asignación fija.}}\\
\rowcolor{gris}
\ch & \textbf{Domingo 4} & libre & \multicolumn{4}{l@{}}{\textit{Comodín. Los patrones 4 y 5 deben quedar cerrados hoy.}}\\
\end{tabular}
\vspace{2mm}

\noindent\begin{tabular}{@{}p{2.0cm}p{22.8cm}@{}}
\textbf{\small Notas} & \rule{\linewidth}{0.4pt}\\[5mm]
 & \rule{\linewidth}{0.4pt}\\[5mm]
 & \rule{\linewidth}{0.4pt}\\[5mm]
 & \rule{\linewidth}{0.4pt}\\[5mm]
 & \rule{\linewidth}{0.4pt}\\
\end{tabular}
\newpage
\seccion{S7 \quad\textnormal{\normalsize 5 de octubre -- 11 de octubre}}
\textbf{Foco:} Relatividad especial\\[0.5mm]
\textbf{Por qué:} Cierra E\&M y prepara las preguntas 11 y 22. Punto medio del plan: momento de reajustar.\\[0.5mm]
\textbf{Entregable de la semana:} 04\_ESTUDIO/formularios/relatividad.md. \textbf{Revisa métricas y reajusta el plan si vas atrasado.}
\par\vspace{2.5mm}
\noindent\begin{tabular}{@{}p{0.45cm}p{2.5cm}p{1.05cm}|p{0.45cm}p{9.75cm}|p{0.45cm}p{9.75cm}@{}}
\rowcolor{cab}
\multicolumn{3}{@{}l}{\color{white}\textbf{\ Día}} & \multicolumn{2}{l}{\color{white}\textbf{MÍNIMA --- 1 hora}} & \multicolumn{2}{l}{\color{white}\textbf{EXTENDIDA --- 2 horas (la mínima + esto)}}\\
\ch & \textbf{Lunes 5} & 16:00 & \ch & Griffiths §12.1 (p.\,477): postulados, geometría, transformaciones de Lorentz (§12.1.3, p.\,493). 4 problemas. & \ch & Además: MIT 8.20 pset1 con su solución.\\
\rowcolor{gris2}
\ch & \textbf{Martes 6} & 16:00 & \ch & Dilatación y contracción. Practica la expansión para $\beta\ll1$: $\gamma\approx1+\beta^2/2$ --- es la \textbf{pregunta 11}. 6 problemas. & \ch & Además: MIT 8.20 pset2 con su solución.\\
\ch & \textbf{Miércoles 7} & 16:00 & \ch & Griffiths §12.2 (p.\,507): energía y momento relativistas, $E^2=(pc)^2+(mc^2)^2$. 6 problemas. & \ch & Además: Irodov §1.8 (p.\,67), 6 problemas.\\
\rowcolor{gris2}
\ch & \textbf{Jueves 8} & 19:00 & \ch & Doppler relativista: factor $\sqrt{(1+\beta)/(1-\beta)}$. Es la base de la \textbf{pregunta 22}. Deriva y haz 4 problemas. & \ch & Además: MIT 8.20 midterm1 con su solución.\\
\ch & \textbf{Viernes 9} & 19:00 & \ch & Banco HRW cap.\,37: 25 preguntas a 2 min. Luego preguntas de relatividad de los 4 GRE. & \ch & Además: MIT 8.20 midterm2 con su solución.\\
\rowcolor{libre}
 & \textbf{Sábado 10} & --- & \multicolumn{4}{l@{}}{\textit{Comodín. Sin asignación fija.}}\\
\rowcolor{gris}
\ch & \textbf{Domingo 11} & libre & \multicolumn{4}{l@{}}{\textit{Comodín. \textbf{Corte de control:} corre \texttt{medidor.py reporte} y lee las recomendaciones. Si algún área va bajo 55\,\%, róbale sesiones a óptica (semana 13).}}\\
\end{tabular}
\vspace{2mm}

\noindent\begin{tabular}{@{}p{2.0cm}p{22.8cm}@{}}
\textbf{\small Notas} & \rule{\linewidth}{0.4pt}\\[5mm]
 & \rule{\linewidth}{0.4pt}\\[5mm]
 & \rule{\linewidth}{0.4pt}\\[5mm]
 & \rule{\linewidth}{0.4pt}\\[5mm]
 & \rule{\linewidth}{0.4pt}\\
\end{tabular}
\newpage
\seccion{S8 \quad\textnormal{\normalsize 12 de octubre -- 18 de octubre}}
\textbf{Foco:} Termodinámica\\[0.5mm]
\textbf{Por qué:} Bloque de 20\,\%. Schroeder cubre esta semana y la siguiente completas.\\[0.5mm]
\textbf{Entregable de la semana:} 04\_ESTUDIO/formularios/termodinamica.md. \textbf{Pide ya las 2 cartas de referencia} --- dependen de terceros.
\par\vspace{2.5mm}
\noindent\begin{tabular}{@{}p{0.45cm}p{2.5cm}p{1.05cm}|p{0.45cm}p{9.75cm}|p{0.45cm}p{9.75cm}@{}}
\rowcolor{cab}
\multicolumn{3}{@{}l}{\color{white}\textbf{\ Día}} & \multicolumn{2}{l}{\color{white}\textbf{MÍNIMA --- 1 hora}} & \multicolumn{2}{l}{\color{white}\textbf{EXTENDIDA --- 2 horas (la mínima + esto)}}\\
\rowcolor{libre}
\ch & \textbf{Lunes 12} & --- & \multicolumn{4}{l@{}}{\textit{\textbf{LIBRE} --- día 12 del mes. La sesión de hoy se corre al domingo comodín.}}\\
\rowcolor{gris2}
\ch & \textbf{Martes 13} & 16:00 & \ch & Schroeder §1.4--1.5 (p.\,17--27): calor, trabajo, $W=-\int p\,dV$. Es la \textbf{pregunta 18}. 6 problemas. & \ch & Además: Irodov §2.2 (p.\,78), 6 problemas.\\
\ch & \textbf{Miércoles 14} & 16:00 & \ch & Schroeder §1.6 (p.\,28): capacidades caloríficas y calor latente. Calorimetría con cambios de fase --- \textbf{pregunta 17}. 6 problemas. & \ch & Además: Lim \emph{Thermodynamics}, parte I: 6 problemas.\\
\rowcolor{gris2}
\ch & \textbf{Jueves 15} & 19:00 & \ch & Schroeder cap.\,4 (p.\,122): máquinas térmicas y ciclo de Carnot. 5 problemas. & \ch & Además: cap.\,5 §5.3 (p.\,166), Clausius--Clapeyron.\\
\ch & \textbf{Viernes 16} & 19:00 & \ch & Banco HRW caps.\,18--20: 30 preguntas a 2 min. & \ch & 60 preguntas.\\
\rowcolor{libre}
 & \textbf{Sábado 17} & --- & \multicolumn{4}{l@{}}{\textit{Comodín. Sin asignación fija.}}\\
\rowcolor{gris}
\ch & \textbf{Domingo 18} & libre & \multicolumn{4}{l@{}}{\textit{Comodín. Formulario de termodinámica.}}\\
\end{tabular}
\vspace{2mm}

\noindent\begin{tabular}{@{}p{2.0cm}p{22.8cm}@{}}
\textbf{\small Notas} & \rule{\linewidth}{0.4pt}\\[5mm]
 & \rule{\linewidth}{0.4pt}\\[5mm]
 & \rule{\linewidth}{0.4pt}\\[5mm]
 & \rule{\linewidth}{0.4pt}\\[5mm]
 & \rule{\linewidth}{0.4pt}\\
\end{tabular}
\newpage
\seccion{S9 \quad\textnormal{\normalsize 19 de octubre -- 25 de octubre}}
\textbf{Foco:} Física estadística\\[0.5mm]
\textbf{Por qué:} Tres patrones caros seguidos: entropía combinatoria, contacto térmico y estadística cuántica.\\[0.5mm]
\textbf{Entregable de la semana:} 04\_ESTUDIO/formularios/estadistica.md + 3 tarjetas Anki.
\par\vspace{2.5mm}
\noindent\begin{tabular}{@{}p{0.45cm}p{2.5cm}p{1.05cm}|p{0.45cm}p{9.75cm}|p{0.45cm}p{9.75cm}@{}}
\rowcolor{cab}
\multicolumn{3}{@{}l}{\color{white}\textbf{\ Día}} & \multicolumn{2}{l}{\color{white}\textbf{MÍNIMA --- 1 hora}} & \multicolumn{2}{l}{\color{white}\textbf{EXTENDIDA --- 2 horas (la mínima + esto)}}\\
\ch & \textbf{Lunes 19} & 16:00 & \ch & Schroeder §2.1--2.2 (p.\,49--55): multiplicidad, sistemas de dos estados, sólido de Einstein. 4 problemas. & \ch & Además: §2.4, aproximación de Stirling (p.\,60).\\
\rowcolor{gris2}
\ch & \textbf{Martes 20} & 16:00 & \ch & \textbf{PATRÓN 6 --- Entropía combinatoria.} Schroeder §2.6 (p.\,74): $S=k\ln\Omega$. Es la \textbf{pregunta 19}. Resuelve el caso de las monedas. & \ch & Además: MIT 8.044 ps1 con su solucionario pss1.\\
\ch & \textbf{Miércoles 21} & 16:00 & \ch & \textbf{PATRÓN 7 --- Contacto térmico.} Schroeder §2.3 (p.\,56) y §3.1 (p.\,85): maximizar $\Omega_A\Omega_B$. Es la \textbf{pregunta 20}. & \ch & Además: MIT 8.044 ps2 con pss2.\\
\rowcolor{gris2}
\ch & \textbf{Jueves 22} & 19:00 & \ch & Schroeder cap.\,6 (p.\,220): factor de Boltzmann, función de partición, distribución de Maxwell (§6.4, p.\,242). 5 problemas. & \ch & Además: Lim \emph{Thermodynamics}, parte II: 6 problemas.\\
\ch & \textbf{Viernes 23} & 19:00 & \ch & \textbf{PATRÓN 8 --- Bosones y fermiones.} Schroeder §7.2 (p.\,262) y §7.4, cuerpo negro (p.\,288). Luego preguntas de estadística de los 4 GRE. & \ch & Además: MIT 8.044 exam1\_03 con su solución.\\
\rowcolor{libre}
 & \textbf{Sábado 24} & --- & \multicolumn{4}{l@{}}{\textit{Comodín. Sin asignación fija.}}\\
\rowcolor{gris}
\ch & \textbf{Domingo 25} & libre & \multicolumn{4}{l@{}}{\textit{Comodín. Formulario de estadística y las 3 tarjetas.}}\\
\end{tabular}
\vspace{2mm}

\noindent\begin{tabular}{@{}p{2.0cm}p{22.8cm}@{}}
\textbf{\small Notas} & \rule{\linewidth}{0.4pt}\\[5mm]
 & \rule{\linewidth}{0.4pt}\\[5mm]
 & \rule{\linewidth}{0.4pt}\\[5mm]
 & \rule{\linewidth}{0.4pt}\\[5mm]
 & \rule{\linewidth}{0.4pt}\\
\end{tabular}
\newpage
\seccion{S10 \quad\textnormal{\normalsize 26 de octubre -- 1 de noviembre}}
\textbf{Foco:} Física moderna y cuántica I\\[0.5mm]
\textbf{Por qué:} Arranca el bloque más caro: todas las preguntas de cuántica son de nivel superior.\\[0.5mm]
\textbf{Entregable de la semana:} 04\_ESTUDIO/formularios/moderna.md. \textbf{Borrador de la carta de intención} (máx.\ 1000 palabras, nombra grupos de investigación).
\par\vspace{2.5mm}
\noindent\begin{tabular}{@{}p{0.45cm}p{2.5cm}p{1.05cm}|p{0.45cm}p{9.75cm}|p{0.45cm}p{9.75cm}@{}}
\rowcolor{cab}
\multicolumn{3}{@{}l}{\color{white}\textbf{\ Día}} & \multicolumn{2}{l}{\color{white}\textbf{MÍNIMA --- 1 hora}} & \multicolumn{2}{l}{\color{white}\textbf{EXTENDIDA --- 2 horas (la mínima + esto)}}\\
\ch & \textbf{Lunes 26} & 16:00 & \ch & Eisberg \& Resnick: radiación de cuerpo negro y efecto fotoeléctrico. 6 problemas. & \ch & Además: efecto Compton, 6 problemas.\\
\rowcolor{gris2}
\ch & \textbf{Martes 27} & 16:00 & \ch & Eisberg: modelo de Bohr, series espectrales, energía de ionización, De Broglie. 6 problemas. & \ch & Además: Lim \emph{Atomic}, sección de física atómica: 5 problemas.\\
\ch & \textbf{Miércoles 28} & 16:00 & \ch & Griffiths QM cap.\,1 y §2.1--2.2: función de onda, Schrödinger, pozo infinito. Memoriza el espectro. 4 problemas. & \ch & Además: Lim \emph{QM} 1001--1071, 5 problemas.\\
\rowcolor{gris2}
\ch & \textbf{Jueves 29} & 19:00 & \ch & \textbf{PATRÓN 9 --- Escalón de potencial.} Griffiths QM cap.\,2. Deriva $R=\left(\frac{k_1-k_2}{k_1+k_2}\right)^2$ y resuelve el caso $E=9V_0/8$ (da $k_1/k_2=3$, luego $R=1/4$) --- es la \textbf{pregunta 23}. & \ch & Además: barrera y tunelamiento, 4 problemas.\\
\ch & \textbf{Viernes 30} & 19:00 & \ch & Banco HRW caps.\,38--40: 25 preguntas a 2 min. Luego preguntas de moderna de los 4 GRE. & \ch & Además: MIT 8.04 ps1 a ps3.\\
\rowcolor{libre}
 & \textbf{Sábado 31} & --- & \multicolumn{4}{l@{}}{\textit{Comodín. Sin asignación fija.}}\\
\rowcolor{gris}
\ch & \textbf{Domingo 1} & libre & \multicolumn{4}{l@{}}{\textit{Comodín. Formulario de moderna. Avanza la carta de intención.}}\\
\end{tabular}
\vspace{2mm}

\noindent\begin{tabular}{@{}p{2.0cm}p{22.8cm}@{}}
\textbf{\small Notas} & \rule{\linewidth}{0.4pt}\\[5mm]
 & \rule{\linewidth}{0.4pt}\\[5mm]
 & \rule{\linewidth}{0.4pt}\\[5mm]
 & \rule{\linewidth}{0.4pt}\\[5mm]
 & \rule{\linewidth}{0.4pt}\\
\end{tabular}
\newpage
\seccion{S11 \quad\textnormal{\normalsize 2 de noviembre -- 8 de noviembre}}
\textbf{Foco:} Cuántica II\\[0.5mm]
\textbf{Por qué:} Cierra los 10 patrones. Aquí están las preguntas 21 y 24.\\[0.5mm]
\textbf{Entregable de la semana:} 04\_ESTUDIO/formularios/cuantica.md + 4 tarjetas Anki. \textbf{Los 10 patrones completos.}
\par\vspace{2.5mm}
\noindent\begin{tabular}{@{}p{0.45cm}p{2.5cm}p{1.05cm}|p{0.45cm}p{9.75cm}|p{0.45cm}p{9.75cm}@{}}
\rowcolor{cab}
\multicolumn{3}{@{}l}{\color{white}\textbf{\ Día}} & \multicolumn{2}{l}{\color{white}\textbf{MÍNIMA --- 1 hora}} & \multicolumn{2}{l}{\color{white}\textbf{EXTENDIDA --- 2 horas (la mínima + esto)}}\\
\ch & \textbf{Lunes 2} & 16:00 & \ch & \textbf{PATRÓN 10 --- Oscilador armónico.} Griffiths QM §2.3.1, método algebraico: $a$, $a^\dagger$, espectro. Deriva $\langle x\rangle(t)$. & \ch & Además: evolución de $\langle(\Delta x)^2\rangle(t)$ --- es la \textbf{pregunta 24}. Resuélvela entera.\\
\rowcolor{gris2}
\ch & \textbf{Martes 3} & 16:00 & \ch & Griffiths QM cap.\,3: formalismo, operadores hermíticos, notación de Dirac, incertidumbre. 4 problemas. & \ch & Además: MIT 8.05 Chap\_01 a Chap\_03.\\
\ch & \textbf{Miércoles 4} & 16:00 & \ch & Griffiths QM cap.\,4: Schrödinger 3D, átomo de hidrógeno, momento angular. Autovalores de $L^2$ y $L_z$. 4 problemas. & \ch & Además: Lim \emph{QM} 3001--3048, 5 problemas.\\
\rowcolor{gris2}
\ch & \textbf{Jueves 5} & 19:00 & \ch & Griffiths QM §5.1, partículas idénticas. Bosones vs.\ fermiones en caja 2D --- es la \textbf{pregunta 21}. Resuélvela. & \ch & Además: Lim \emph{QM} 7001--7037, 4 problemas.\\
\ch & \textbf{Viernes 6} & 19:00 & \ch & Preguntas de cuántica de los 4 GRE, cronometradas. & \ch & Además: Cahn \& Nadgorny parte 2, cap.\,5: problemas 5.1 a 5.21.\\
\rowcolor{libre}
 & \textbf{Sábado 7} & --- & \multicolumn{4}{l@{}}{\textit{Comodín. Sin asignación fija.}}\\
\rowcolor{gris}
\ch & \textbf{Domingo 8} & libre & \multicolumn{4}{l@{}}{\textit{Comodín. Repasa los 10 patrones de corrido. Si alguno falla, hoy se arregla.}}\\
\end{tabular}
\vspace{2mm}

\noindent\begin{tabular}{@{}p{2.0cm}p{22.8cm}@{}}
\textbf{\small Notas} & \rule{\linewidth}{0.4pt}\\[5mm]
 & \rule{\linewidth}{0.4pt}\\[5mm]
 & \rule{\linewidth}{0.4pt}\\[5mm]
 & \rule{\linewidth}{0.4pt}\\[5mm]
 & \rule{\linewidth}{0.4pt}\\
\end{tabular}
\newpage
\seccion{S12 \quad\textnormal{\normalsize 9 de noviembre -- 15 de noviembre}}
\textbf{Foco:} Óptica express y primeros simulacros\\[0.5mm]
\textbf{Por qué:} Óptica en dos sesiones (pesa 0--4\,\%) y arranque de simulacros completos.\\[0.5mm]
\textbf{Entregable de la semana:} 04\_ESTUDIO/formularios/optica.md (media página basta) + 2 simulacros registrados.
\par\vspace{2.5mm}
\noindent\begin{tabular}{@{}p{0.45cm}p{2.5cm}p{1.05cm}|p{0.45cm}p{9.75cm}|p{0.45cm}p{9.75cm}@{}}
\rowcolor{cab}
\multicolumn{3}{@{}l}{\color{white}\textbf{\ Día}} & \multicolumn{2}{l}{\color{white}\textbf{MÍNIMA --- 1 hora}} & \multicolumn{2}{l}{\color{white}\textbf{EXTENDIDA --- 2 horas (la mínima + esto)}}\\
\ch & \textbf{Lunes 9} & 16:00 & \ch & Óptica geométrica: Sears Vol.\,2. Snell, reflexión total interna, lentes delgadas, aumento. Formulario de media página. & \ch & Además: Delft \emph{BSc Optics} cap.\,2 (p.\,39), §2.5 óptica gaussiana (p.\,46).\\
\rowcolor{gris2}
\ch & \textbf{Martes 10} & 16:00 & \ch & Interferencia y difracción: Young, red de difracción, criterio de Rayleigh, Malus, Brewster. Añade al formulario. 6 problemas. & \ch & Además: Irodov §5.2--5.4 (p.\,210--226), 6 problemas.\\
\ch & \textbf{Miércoles 11} & 16:00 & \ch & \textbf{SIMULACRO 1:} GR8677, primera mitad. 85 min cronometrados, condiciones reales. & \ch & El examen completo: 170 min seguidos.\\
\rowcolor{libre}
\ch & \textbf{Jueves 12} & --- & \multicolumn{4}{l@{}}{\textit{\textbf{LIBRE} --- día 12 del mes. La sesión de hoy se corre al domingo comodín.}}\\
\ch & \textbf{Viernes 13} & 19:00 & \ch & \textbf{SIMULACRO 2:} REA Test 1 (p.\,79), primera mitad. 60 min. & \ch & Completo, más la corrección con las explicaciones detalladas (p.\,110).\\
\rowcolor{libre}
 & \textbf{Sábado 14} & --- & \multicolumn{4}{l@{}}{\textit{Comodín. Sin asignación fija.}}\\
\rowcolor{gris}
\ch & \textbf{Domingo 15} & libre & \multicolumn{4}{l@{}}{\textit{Comodín. \textbf{Aviso:} la inscripción online cierra el miércoles 18 de noviembre a las 5:00 p.m.}}\\
\end{tabular}
\vspace{2mm}

\noindent\begin{tabular}{@{}p{2.0cm}p{22.8cm}@{}}
\textbf{\small Notas} & \rule{\linewidth}{0.4pt}\\[5mm]
 & \rule{\linewidth}{0.4pt}\\[5mm]
 & \rule{\linewidth}{0.4pt}\\[5mm]
 & \rule{\linewidth}{0.4pt}\\[5mm]
 & \rule{\linewidth}{0.4pt}\\
\end{tabular}
\newpage
\seccion{S13 \quad\textnormal{\normalsize 16 de noviembre -- 22 de noviembre}}
\textbf{Foco:} Simulacros finales, trámites y repaso\\[0.5mm]
\textbf{Por qué:} Nada nuevo. Solo simulacros, corrección de errores y cerrar la inscripción.\\[0.5mm]
\textbf{Entregable de la semana:} Inscripción y documentos entregados. Los 5 formularios repasados.
\par\vspace{2.5mm}
\noindent\begin{tabular}{@{}p{0.45cm}p{2.5cm}p{1.05cm}|p{0.45cm}p{9.75cm}|p{0.45cm}p{9.75cm}@{}}
\rowcolor{cab}
\multicolumn{3}{@{}l}{\color{white}\textbf{\ Día}} & \multicolumn{2}{l}{\color{white}\textbf{MÍNIMA --- 1 hora}} & \multicolumn{2}{l}{\color{white}\textbf{EXTENDIDA --- 2 horas (la mínima + esto)}}\\
\ch & \textbf{Lunes 16} & 16:00 & \ch & \textbf{SIMULACRO 3:} GR9677, primera mitad. 85 min. & \ch & El examen completo: 170 min.\\
\rowcolor{gris2}
\ch & \textbf{Martes 17} & 16:00 & \ch & Corrige el simulacro 3. Repasa los 10 patrones en voz alta, sin mirar. & \ch & Además: REA Test 2 (p.\,153), primera mitad.\\
\ch & \textbf{Miércoles 18} & 16:00 & \ch & \textcolor{alerta}{\textbf{HOY CIERRA LA INSCRIPCIÓN ONLINE, 5:00 p.m. Hazlo antes de estudiar.}} Luego: repaso de errores.md, rehaz los 10 fallos más repetidos. & \ch & Además: REA Test 3 (p.\,233), primera mitad.\\
\rowcolor{gris2}
\ch & \textbf{Jueves 19} & 19:00 & \ch & \textcolor{alerta}{\textbf{HOY CIERRA LA CARGA DE DOCUMENTOS, 11:30 p.m.}} Luego: repasa los 5 formularios. Nada nuevo. & \ch & Además: pasa todas las tarjetas Anki.\\
\ch & \textbf{Viernes 20} & 19:00 & \ch & Repaso ligero: los 10 patrones y la tabla de constantes (REA p.\,76). \textbf{Máximo 1 hora. No estudies nada nuevo.} & \ch & \textbf{Hoy no hay versión extendida.} Descansa: rinde más que estudiar.\\
\rowcolor{libre}
 & \textbf{Sábado 21} & --- & \multicolumn{4}{l@{}}{\textit{Comodín. Sin asignación fija.}}\\
\rowcolor{gris}
\ch & \textbf{Domingo 22} & libre & \multicolumn{4}{l@{}}{\textit{Descanso. Prepara documento de identidad, lápiz, borrador y reloj. Duerme temprano: el examen es mañana.}}\\
\end{tabular}
\vspace{2mm}

\noindent\begin{tabular}{@{}p{2.0cm}p{22.8cm}@{}}
\textbf{\small Notas} & \rule{\linewidth}{0.4pt}\\[5mm]
 & \rule{\linewidth}{0.4pt}\\[5mm]
 & \rule{\linewidth}{0.4pt}\\[5mm]
 & \rule{\linewidth}{0.4pt}\\[5mm]
 & \rule{\linewidth}{0.4pt}\\
\end{tabular}
\newpage

\begin{center}{\color{alerta}\Huge\textbf{LUNES 23 DE NOVIEMBRE DE 2026 --- EXAMEN}}\\[2mm]
{\large Presencial. Hora y salón se informan el día anterior. 25 preguntas, 3 horas.}\end{center}
\vspace{1mm}
\footnotesize

\begin{minipage}[t]{0.48\textwidth}
\seccion{Los 10 patrones --- lista de verificación}
Si el 22 de noviembre solo alcanzas a repasar diez cosas, son estas. Marca cuando puedas explicarla sin mirar:
\begin{enumerate}[leftmargin=6mm,itemsep=1.2mm]
\item \ch\ $T(E)=dA/dE$ en el espacio de fase
\item \ch\ Corchetes de Poisson: $\{L_i,p_j\}=\varepsilon_{ijk}p_k$
\item \ch\ Scattering esfera dura: $b=R\cos(\theta/2)$
\item \ch\ Condición de gauge de Lorenz y de Coulomb
\item \ch\ Esfera conductora en campo uniforme: $(Ar+B/r^2)\cos\theta$
\item \ch\ $S=k\ln\Omega$ y entropía combinatoria
\item \ch\ Equilibrio térmico maximizando $\Omega_A\Omega_B$
\item \ch\ Bosones vs.\ fermiones en caja: quién comparte estado
\item \ch\ Coeficiente $R$ del escalón de potencial
\item \ch\ Oscilador armónico con $a$, $a^\dagger$ y evolución temporal
\end{enumerate}

\seccion{Las 4 técnicas de descarte}
Con 7,2 min por pregunta, muchas veces descartar es más rápido que resolver:
\begin{enumerate}[leftmargin=6mm,itemsep=0.8mm]
\item \textbf{Dimensional:} ¿las unidades de cada opción cierran?
\item \textbf{Casos límite:} $x\to0$, $x\to\infty$, $m_1=m_2$, sin fricción, $v\ll c$, $T\to0$.
\item \textbf{Signos y simetrías:} ¿el campo apunta a donde debe? ¿la energía crece o decrece?
\item \textbf{Orden de magnitud:} descarta lo que difiere por factores absurdos.
\end{enumerate}
\textbf{Responde las 25.} No hay penalización por error: dejar en blanco solo pierde puntos.
\end{minipage}
\hfill
\begin{minipage}[t]{0.48\textwidth}
\seccion{Fechas administrativas --- no se negocian}
\begin{tabular}{@{}p{3.2cm}p{9cm}@{}}
\textbf{18 de nov, 5:00 p.m.} & \ch\ Cierre de la inscripción online\\
\textbf{19 de nov, 11:30 p.m.} & \ch\ Cierre de carga de documentos\\
\textbf{23 de nov} & \ch\ Examen de admisión\\
\textbf{1--2 de dic} & Entrevistas\\
\textbf{2 de dic} & Publicación de admitidos\\
\end{tabular}
\vspace{2mm}

\textbf{Documentos (PDF, máximo 1 MB cada uno):}
\begin{itemize}[leftmargin=6mm,itemsep=0.4mm]
\item \ch\ Diploma y/o acta de grado, con firma digital validable en línea
\item \ch\ Certificado oficial de calificaciones de pregrado
\item \ch\ Hoja de vida
\item \ch\ Carta de intención (máx.\ 1000 palabras; debe nombrar grupos de investigación)
\item \ch\ Dos referencias académicas, desde correos institucionales
\item \ch\ Carta de solicitud de apoyo financiero (opcional)
\end{itemize}
\textbf{Tenlos listos antes del 31 de octubre.} Las referencias dependen de terceros: pídelas en la semana 9.
\vspace{2mm}

\seccion{Entregables por semana}
\begin{tabular}{@{}p{1.3cm}p{11cm}@{}}
S1 & \ch\ Diagnóstico corregido y registrado\\
S2 & \ch\ Formulario mecánica I\\
S3 & \ch\ Formulario mecánica II\\
S4 & \ch\ Formulario mecánica analítica + patrones 1--3\\
S5 & \ch\ Formulario electrostática\\
S6 & \ch\ Formulario magnetostática e inducción\\
S7 & \ch\ Formulario Maxwell/ondas + patrones 4--5\\
S8 & \ch\ Formulario relatividad + corte de control\\
S9 & \ch\ Formulario termodinámica + cartas pedidas\\
S10 & \ch\ Formulario estadística + patrones 6--8\\
S11 & \ch\ Formulario moderna + patrón 9 + carta de intención\\
S12 & \ch\ Formulario cuántica + patrón 10\\
S13 & \ch\ Formulario óptica + simulacros 1 y 2\\
S14 & \ch\ Simulacros 3 y 4 + inscripción cerrada\\
\end{tabular}
\end{minipage}

\end{document}
